\documentclass[10pt,a4paper]{amsart}
\usepackage[margin=19mm]{geometry}
\usepackage{amsmath,amssymb,mathtools}
\usepackage[scr=rsfs,cal=euler]{mathalfa}
\usepackage{xfrac}
\usepackage[unicode,hidelinks,bookmarks=false]{hyperref}
\newcommand{\ie}{i.e.\ }
\newcommand{\cf}{cf.\ }
\newcommand{\resp}{respectively\ }
\newcommand{\Subsection}{\subsection}
\providecommand{\nobreakdash}{\nobreak\hbox{-}}
\setlength{\emergencystretch}{2em}
\multlinegap0pt
\usepackage{bm}
\usepackage{bbm}
\usepackage{dsfont}
\usepackage{xspace}
\usepackage{cancel}
\usepackage{mathtools}
\usepackage{amsthm}
\makeatletter
\@ifundefined{lema}{\newtheorem{lema}{Lemma}}{}
\makeatother
\title[Controllability of a system of non-autonomous degenerate cou]{Controllability of a system of non-autonomous degenerate coupled parabolic equations}
\author{Alfredo S. Gamboa, Juan Limaco, Luis P. Yapu}
\date{Source: arXiv:2508.17546v1 (CC BY 4.0)}
\begin{document}
\maketitle

\noindent\emph{Source and licence.} Controllability of a system of non-autonomous degenerate coupled parabolic equations. Version arXiv:2508.17546v1 (CC BY 4.0), distributed under the Creative Commons Attribution 4.0 International licence, as recorded in the arXiv licence metadata for this exact version and shown on the abstract page https://arxiv.org/abs/2508.17546v1. Full rights notice: licenses/arxiv-2508.17546-CC-BY-4.0.txt.\par\smallskip

\noindent\textit{Editorial scope.} Counted unit: the Lema whose source label is \texttt{lemaA8} in the article (source0.tex lines 765--776, source label \texttt{lemaA8}), delivered together with the article's own complete original proof of it (source0.tex lines 777--867, one \texttt{proof} environment of 91 lines). No mathematics is written, repaired or re-derived: statement and proof are the article's own text line for line. Required context retained uncounted: eq:int\_I4 (lines 686--689). Every definition, display and label of those lines is retained unchanged. Omitted from this package and not redistributed: 1--685, 690--764, 868--1989. Nothing inside a retained excerpt is omitted or altered: every retained body is delivered exactly as the article's lines are written, so no omission receipt of any kind accompanies this package. Citations: the retained excerpts carry no citation command at all, so no bibliography entry of the article is transcribed and no reference is redistributed. Reference adaptation: none was needed and none was made; every reference inside the delivered unit resolves inside it, so no unresolved reference or citation is produced. Printed slips kept literal: no printed word, symbol, hypothesis or number of the counted statement or of its proof was altered. Layout and class: the article's own class is \texttt{osajnl} with packages including amsthm, comment; the wrapper is the lane's pinned \texttt{amsart} editorial wrapper at 10pt on A4, which restates the theorem-like environments the retained text needs and adds the article's own macro definitions that the retained text needs (none), with extra wrapper packages (bm, bbm, dsfont, xspace, cancel, mathtools). The delivered numbering is the unit's own. Assets: no retained span contains a figure environment, a graphics-inclusion command, a PDF-page inclusion, a table or a TikZ picture; no image file is copied into this package, the package declares no asset dependency and no image file is redistributed with it. Licence: version arXiv:2508.17546v1 of this article is distributed under the Creative Commons Attribution 4.0 International licence, as recorded in the arXiv licence metadata for this exact version and shown on the abstract page https://arxiv.org/abs/2508.17546v1. A full rights notice is delivered at licenses/arxiv-2508.17546-CC-BY-4.0.txt. Characters: every retained excerpt is pure ASCII, so the delivered engine, XeLaTeX with its default Latin Modern text font, reports no missing character.
\begin{equation}
\label{eq:int_I4}
I_4=s\int_{Q}ab^2(a\varphi_x)_{xx} w_x w+2s\int_{Q}ab^2(a\varphi_x )_xw_x^2-s\int_{Q}ab^2 a_x \varphi_xw_{x}^2-s\int_{0}^{T}\left[a^2b^2\varphi_x w_x^2\right]^{x=1}_{x=0}.
\end{equation}

\begin{lema}[Estimate for $I_4$ of \eqref{eq:int_I4}]
\label{lemaA8}
	\begin{equation*}
    \begin{split}
		s \int_Q b^2 a (a \varphi_x)_{xx} w_x w \geq
		&- C s^2\lambda^3 \int_0^T \int_0^{\alpha'} b^2 \frac{x^2}{a}  \sigma^3 w^2 - C \lambda \int_0^T \int_0^{\alpha'} b^2 a \sigma w_x^2
		- Cs^2\lambda^3 \int_0^T\int_{\omega'} b^2 \sigma^3 w^2 \\
		&- C\lambda \int_0^T\int_{\omega'} b^2 \sigma w_x^2
		-Cs^2\lambda^4 \int_Q b^2 a^2 |\psi_x|^4 \sigma^3 w^2 - C\lambda^2 \int_Q b^2 a^2|\psi_x|^2\sigma w_x^2.
    \end{split}
	\end{equation*}
\end{lema}

\begin{proof}
    From the definition of $\varphi$, we have
    $$\varphi_x=\lambda \psi_x \sigma, \ \ \ \ \ \ \varphi_{xx}=\lambda \psi_{xx}\sigma +\lambda^2\psi_{xx}\psi_x \sigma, \ \ \ \ \ \ \varphi_{xxx}=\lambda \psi_{xxx}\sigma +3\lambda^2 \psi_{xx}\psi_x \sigma+\lambda^3\psi_x^3\sigma,$$
    %$$(a\varphi_x)_{xx}=a''\varphi_x+2a'\varphi_{xx}+a\varphi_{xxx}, \ \ \ \ $$
    $$ (a\varphi_x)_{xx}=\left[ \lambda (a\psi_x)_{xx}+2\lambda^2 \psi_x (a\psi_x)_x+\lambda^2a\psi_x\psi_{xx}+\lambda^3 a\psi_x^3   \right]\sigma.$$
    
    Thus, 
	\begin{align*}
		s\int_Q b^2 a (a \varphi_x)_{xx} w_x w=&s\lambda\int_Q ab^2 (a\psi_x)_{xx} \sigma w_x w+2s\lambda^2 \int_Q ab^2 \psi_x (a\psi_x)_x \sigma w_x w\\
		&+s\lambda^2\int_Q a^2b^2\psi_x\psi_{xx} \sigma w_x w
		+s\lambda^3\int_Q a^2b^2\psi_x^3 \sigma w_x w.
	\end{align*}
    
	From the definition of $\psi$ in $(0,\alpha')$ and $(\beta',1)$, we have that $\psi_x=\pm\frac{x}{a}$. With this we get $(a \varphi_x)_{xx}=0$  in those domains. Thus,
		
%	\noindent The first integral in the inequality implies that, 
%	\begin{multline*}
%		s\lambda\int_Q ab^2 (a\psi_x)_{xx} \sigma w_x w=s\lambda\left[\int_0^T \int_0^{\alpha'} ab^2 (a\psi_x)_{xx} \sigma w_x w +  \int_0^T\int_{\omega'} ab^2 (a\psi_x)_{xx} \sigma w_x w \right.\\
%		\left.+\int_0^T \int_{\beta'}^1 ab^2 (a\psi_x)_{xx} \sigma w_x w\right]
%	\end{multline*}
	$$s\lambda\int_Q ab^2 (a\psi_x)_{xx} \sigma w_x w=s\lambda\int_0^T\int_{\omega'} ab^2 (a\psi_x)_{xx} \sigma w_x w $$
%	since $(a\psi_x)_{xx}=0$ in  $\omega'$ in $(\beta', 1).$\\	

	Then, returning to the expression we want to estimate, 
	\begin{align*}
		s\int_Q b^2 a (a \varphi_x)_{xx} w_x w=&s\lambda\int_0^T\int_{\omega'} ab^2 (a\psi_x)_{xx} \sigma w_x w +2s\lambda^2 \int_Q ab^2 \psi_x (a\psi_x)_x \sigma w_x w\\
		&+s\lambda^2\int_Q a^2b^2\psi_x\psi_{xx} \sigma w_x w
		+s\lambda^3\int_Q a^2b^2\psi_x^3 \sigma w_x w
		=J_1+J_2+J_3+J_4.
	\end{align*}
    
    We estimate each integral of the last sum. Taking into account that $\sigma\leq C\sigma^2$, we have
    $$\left|J_1 \right|\leq s\lambda\int_0^T\int_{\omega'} b^2 |a(a\psi_x)_{xx}| \sigma |w_x| |w|\leq Cs\lambda\int_0^T\int_{\omega'} b^2  \sigma^2 |w_x| |w|$$
    Then,
    $$
    \left|J_1 \right|\leq Cs\lambda\int_0^T\int_{\omega'} b \sigma^{3/2}|w| b\sigma^{1/2} |w_x| \leq Cs\lambda\int_0^T\int_{\omega'}b^2 \sigma^3|w|^2 + Cs\lambda\int_0^T\int_{\omega'}b^2 \sigma|w_x|^2.
    $$
	
    \noindent {\bf Estimates for $J_2$:}
    \begin{multline*}
    \left|J_2 \right|\leq s\lambda^2\left[\int_0^T \int_0^{\alpha'} ab^2 \psi_x (a\psi_x)_x \sigma |w_x| |w| +  \int_0^T\int_{\omega'} ab^2 \psi_x (a\psi_x)_x \sigma |w_x| |w| + \int_0^T \int_{\beta'}^1 ab^2 \psi_x (a\psi_x)_x \sigma |w_x| |w|\right]
    \end{multline*}
    On the other hand, we get that $\left| a\psi_x (a\psi_x)_x  \right|=\left| a\cdot \frac{x}{a}\left( a\frac{x}{a} \right)_x \right|=|x|$ for $x\in (0,\alpha')$ and similarly for $x\in (\beta',1)$. We also have the following inequalities: $\sigma\leq C\sigma^2$, \ \ $x\leq a^2|\psi_x|^3$ in $x\in (\beta',1)$. Then,
    \begin{equation*}
    \begin{split}
    \left|J_2 \right| &\leq s\lambda^2\left[\int_0^T \int_0^{\alpha'} b^2 x \sigma |w_x| |w| +  \int_0^T\int_{\omega'} ab^2 \psi_x (a\psi_x)_x \sigma |w_x| |w| + \int_0^T \int_{\beta'}^1 b^2 x \sigma |w_x| |w|\right] \\
    &\leq s\lambda^2 C \left[\int_0^T \int_0^{\alpha'} b^2 x \sigma^2 |w_x| |w| +  \int_0^T\int_{\omega'} b^2  \sigma^2 |w_x| |w| + \int_0^T \int_{\beta'}^1 a^2b^2|\psi_x|^3 \sigma^2 |w_x| |w|\right] \\
    &\leq C \int_0^T \int_0^{\alpha'} \left( \frac{x}{\sqrt{a}}b\sigma^{3/2}\lambda^{3/2}s|w|\right)\left( \sqrt{a}b\sigma^{1/2}\lambda^{1/2}|w_x| \right) \\
    &+C\int_0^T\int_{\omega'} \left( b\sigma^{3/2}\lambda^{3/2}s|w|\right)\left( b\sigma^{1/2}\lambda^{1/2}|w_x| \right)
		+C\int_0^T \int_{\beta'}^1\left( b\sigma^{3/2}sa\lambda^2|\psi_x|^2|w|\right)\left( b\sigma^{1/2}a|\psi_x||w_x| \right)
    \end{split}    
    \end{equation*}
    Thus,
    \begin{equation*}
    \begin{split}
        \left|J_2 \right| \leq &Cs^2\lambda^3\int_0^T \int_0^{\alpha'}\frac{x^2}{a}b^2\sigma^3|w|^2+\lambda\int_0^T \int_0^{\alpha'}ab^2\sigma|w_x|^2+Cs^2\lambda^3\int_0^T\int_{\omega'}b^2\sigma^3|w|^2\\
		&+C\lambda \int_0^T\int_{\omega'}b^2\sigma|w_x|^2
		+Cs^2\lambda^4\int_0^T \int_{\beta'}^1a^2b^2\sigma^3|\psi_x|^4|w|^2+C\int_0^T \int_{\beta'}^1a^2b^2|\psi_x|^2|w_x|^2.    
    \end{split}    
    \end{equation*}

    \noindent {\bf Estimates for $J_3$:}
    $$|J_3|\leq s\lambda^2\int_Q a^2b^2|\psi_x\psi_{xx}| \sigma |w_x| |w|$$
    \begin{multline*}
	|J_3|\leq s\lambda^2\int_0^T \int_0^{\alpha'}          a^2b^2|\psi_x\psi_{xx}| \sigma |w_x| |w| +  s\lambda^2\int_0^T\int_{\omega'} a^2b^2|\psi_x\psi_{xx}| \sigma |w_x| |w|
	+ s\lambda^2\int_0^T \int_{\beta'}^1 a^2b^2|\psi_x\psi_{xx}| \sigma |w_x| |w|.
    \end{multline*}
    
    From the definition of $\psi$, one has,
    $\psi_x\psi_{xx}=\pm \frac{x}{a}\cdot \frac{a-xa'}{a^2}=\pm \frac{x}{a^3}(a-xa')$, for $x\in (0,\alpha')$ and for $x\in (\beta',1)$.
    We also know that $xa'-a\leq Ka-a\leq Ca$. Then,
    \begin{multline*}
		|J_3|\leq s\lambda^2\int_0^T \int_0^{\alpha'} a^2b^2\left|  \frac{a-xa'}{a^3} \right| \sigma |w_x| |w| +  s\lambda^2\int_0^T\int_{\omega'} a^2b^2|\psi_x\psi_{xx}| \sigma |w_x| |w|\\
		+ s\lambda^2\int_0^T \int_{\beta'}^1 a^2b^2\left| - \frac{a-xa'}{a^3} \right| \sigma |w_x| |w|
    \end{multline*}
    $$|J_3|\leq Cs\lambda^2\int_0^T \int_0^{\alpha'} b^2x \sigma^2 |w_x| |w| +  Cs\lambda^2\int_0^T\int_{\omega'} b^2 \sigma^2 |w_x| |w| +C s\lambda^2\int_0^T \int_{\beta'}^1 b^2x \sigma^2 |w_x| |w|,$$
    and we can follow the same estimate as for $J_2$. 
%    The procedure is the same as for $J_2$, that is,
%    \begin{multline*}
%		\left|J_3 \right|\leq Cs^2\lambda^3\int_0^T \int_0^{\alpha'}\frac{x^2}{a}b^2\sigma^3|w|^2+\lambda\int_0^T \int_0^{\alpha'}ab^2\sigma|w_x|^2+Cs^2\lambda^3\int_0^T\int_{\omega'}b^2\sigma^3|w|^2+C\lambda \int_0^T\int_{\omega'}b^2\sigma|w_x|^2\\
%		+Cs^2\lambda^4\int_0^T \int_{\beta'}^1a^2b^2\sigma^3|\psi_x|^4|w|^2+C\int_0^T \int_{\beta'}^1a^2b^2|\psi_x|^2|w_x|^2
%    \end{multline*}
 
    \noindent {\bf Estimates for $J_4:$}
    \begin{equation*}
    \begin{split}
    |J_4| &\leq s\lambda^3\int_Q a^2b^2|\psi_x|^3 \sigma |w_x|  |w|\leq \int_{Q} \left( sab\lambda^{2}|\psi_x|^2\sigma^{3/2}|w|\right)\left( ab\lambda|\psi_x|\sigma^{1/2}|w_x| \right) \\
    &\leq s^2\lambda^4\int_Qa^2b^2|\psi_x|^4 \sigma^3 |w|^2+C\lambda^2\int_{Q} a^2b^2|\psi_x|^2 \sigma |w_x|^2
    \end{split}    
    \end{equation*}
\end{proof}

\end{document}
